% ===========================================================================
\section{Phase 3: Neural Heuristic Model}
\label{sec:phase3_model}
% ===========================================================================

\subsection{Problem formulation}
A$^\ast$ expands nodes in order of
\begin{equation}
  f(n) = g(n) + h(n),
\end{equation}
where $g(n)$ is the exact cost from the start to node $n$ and $h(n)$ estimates
the remaining cost to the goal. On a four-connected grid with unit move cost,
the Manhattan distance
\begin{equation}
  d_M(n) = |r_n-r_{\text{goal}}| + |c_n-c_{\text{goal}}|
\end{equation}
is admissible: every move changes the Manhattan distance by at most one, so no
path can be shorter than $d_M(n)$. It is, however, blind to obstacles. The goal
of this phase is to learn a function $h_\theta(n)$ that approximates the exact
cost-to-goal $h^\ast(n)$ and therefore prunes more of the search space than
Manhattan distance while remaining cheap to evaluate inside A$^\ast$.

We formulate the task as supervised regression over individual grid cells. For
a maze of side length $N$, the regression target is the normalised cost
\begin{equation}
  y(n) = \frac{h^\ast(n)}{N},
\end{equation}
computed for every reachable free cell by a single Dijkstra search from the
goal. Normalising by $N$ keeps the target scale comparable across the five grid
sizes while preserving the relative difficulty of detours.

\subsection{Framework choice}
The model is implemented in PyTorch and trained on CPU. PyTorch provides the
automatic differentiation, checkpointing, and batched inference required for
this project without requiring a GPU. The model is deliberately small: it must
be evaluated for every cell of a maze at search time, so inference cost matters
as much as validation accuracy. A feed-forward MLP is sufficient because the
input is a fixed-length per-cell feature vector, not a variable-length sequence
or image that would justify a convolutional or recurrent architecture.

\subsection{Architecture}
Figure~\ref{fig:architecture} shows the architecture. A shared three-layer
trunk maps the 29-dimensional input to a 128-dimensional representation. Two
heads then branch from this representation: a linear regression head that
predicts the standardised cost-to-goal, and a 10-way classification head that
predicts a distance bin. The auxiliary bin head gives the model a coarse,
ordinal view of the cost and supplies the precision, recall, F1, and AUC-ROC
metrics required by the rubric, while the regression head supplies the scalar
value used by A$^\ast$.

\begin{figure}[H]
  \centering
  \includegraphics[width=\textwidth]{architecture.png}
  \caption{Architecture of the neural heuristic. The shared trunk processes a
  29-dimensional per-cell feature vector. The regression head produces the
  scalar cost used by A$^\ast$; the auxiliary head predicts a distance bin.}
  \label{fig:architecture}
\end{figure}

\begin{table}[H]
  \centering
  \caption{Layer configuration and parameter counts. All linear layers are
  followed by ReLU activations; dropout is applied only in the trunk.}
  \label{tab:model_layers}
  \begin{tabular}{llrr}
    \toprule
    Component & Configuration & Output shape & Parameters \\
    \midrule
    Input               & 29 features                       & $(B, 29)$  & -- \\
    Trunk layer 1       & Linear $29\to256$, ReLU, dropout  & $(B, 256)$ & 7{,}680 \\
    Trunk layer 2       & Linear $256\to256$, ReLU, dropout & $(B, 256)$ & 65{,}792 \\
    Trunk layer 3       & Linear $256\to128$, ReLU          & $(B, 128)$ & 32{,}896 \\
    Regression head     & Linear $128\to1$                  & $(B, 1)$   & 129 \\
    Distance-bin head   & Linear $128\to10$                 & $(B, 10)$  & 1{,}290 \\
    \midrule
    \textbf{Total}      &                                   &            & \textbf{107{,}787} \\
    \bottomrule
  \end{tabular}
\end{table}

The forward pass returns a tuple \texttt{(y\_pred, bin\_logits)}. The regression
output is a single scalar per cell, and the bin head returns ten logits. Dropout
with probability $0.1$ is active during training and disabled at inference, so
predictions are deterministic for a fixed checkpoint.

\subsection{Feature representation}
Each cell is represented by the same 29 features used by the data pipeline:
\begin{enumerate}
  \item the normalised coordinates $(x/N,\,y/N)$ (2 features);
  \item the goal-relative offsets $(\Delta x/N,\,\Delta y/N)$ (2 features);
  \item a $5\times5$ binary wall window centred on the cell, flattened to 25
  features and padded with wall values outside the grid.
\end{enumerate}
The wall window gives the network local obstacle context, the goal offsets tell
it where the target lies, and the coordinates give it global position. The same
builder is used both when generating training data and when constructing
search-time inputs, which removes the risk of train--inference feature skew.
The centre of the wall window is constant for emitted samples because only free
cells are labelled; it is retained for a uniform representation.

\subsection{Training objective and optimisation}
Targets are standardised using the training-set mean $\mu$ and standard
deviation $\sigma$:
\begin{equation}
  y_{\text{std}} = \frac{y-\mu}{\sigma}.
\end{equation}
The regression loss is SmoothL1 with $\beta=0.1$, chosen for robustness to the
heavy tail of long detours. The full training objective is
\begin{equation}
  \mathcal{L}
  = \operatorname{SmoothL1}_\beta(\hat{y}_{\text{std}}, y_{\text{std}})
  + \lambda_{\text{aux}}\,\operatorname{CE}(\text{bin\_logits}, b),
\end{equation}
where $b$ is the distance bin of the target and $\lambda_{\text{aux}}$ balances
the auxiliary classification loss. An optional asymmetric variant adds a
one-sided penalty on overestimation:
\begin{equation}
  \mathcal{L}_{\text{asym}}
  = \mathcal{L}
  + \lambda_{\text{over}}\,\mathbb{E}\!\left[
      \max(0,\ \hat{y}_{\text{std}} - y_{\text{std}} - m)^2
    \right],
\end{equation}
with margin $m$. Overestimation is penalised because it can make A$^\ast$
inadmissible and return a suboptimal path; underestimation only costs extra
node expansions.

Optimisation uses AdamW with a \texttt{ReduceLROnPlateau} scheduler driven by
validation MAE, gradient clipping at $1.0$, and early stopping with patience
five epochs. Dropout in the trunk provides additional regularisation. The best
checkpoint stores the model state, model configuration, target mean and
standard deviation, bin edges, training configuration, git commit, and seed, so
that search-time inference reproduces the exact training-time scaling.

\subsection{Target range and output activation}
The target is not bounded by one. A path may have to detour around obstacles,
so $h^\ast(n)$ can exceed the grid side length $N$; in the processed data the
maximum observed normalised target is approximately $3.18$ (318 cells in a
$100\times100$ maze). The earlier description of the target as lying in
$[0,1]$ is therefore incorrect. The regression head consequently uses a
\emph{linear} output with no sigmoid or other bounded activation. A sigmoid
would saturate exactly on the difficult long-detour cases and would be unable
to represent costs above $N$.

At inference, the prediction is converted back to grid-cell units:
\begin{equation}
  \hat{h}(n) = \max\!\left(0,\ \hat{y}_{\text{std}}\,\sigma + \mu\right) N.
\end{equation}
The clamp enforces non-negativity without affecting valid targets. Empirically,
the training split has $\mu\approx0.738$ and $\sigma\approx0.368$; these values
are recomputed from the actual training run and stored in the checkpoint rather
than hard-coded.

\subsection{Admissibility and search-time safety}
A learned heuristic can overestimate $h^\ast$ on some cells. Because Manhattan
distance is a lower bound on the true cost in a four-connected maze, the
following clipping mode is explicitly safe:
\begin{equation}
  h_{\text{clip}}(n) = \min\!\left(\alpha\,h_\theta(n),\ d_M(n)\right),
  \qquad 0 \le \alpha \le 1.
\end{equation}
Since $d_M(n)\le h^\ast(n)$ for every cell, $h_{\text{clip}}(n)\le h^\ast(n)$,
so this mode is admissible regardless of the network's prediction. The learned
signal can still improve on Manhattan whenever $\alpha h_\theta(n)<d_M(n)$;
otherwise the heuristic falls back smoothly to the classical bound.

Table~\ref{tab:heuristic_modes} summarises the three modes exposed by the
wrapper. The raw mode is the fastest and most aggressive; the scaled mode
shrinks predictions by $\alpha$ to reduce overestimation; the clipped mode
provides the admissibility guarantee above. Cells that are walls, outside the
grid, or have non-finite predictions fall back to Manhattan distance, which
keeps the search well-defined on the complete grid.

\begin{table}[H]
  \centering
  \caption{Search-time heuristic modes. The clipped mode is admissible by
  construction; the other modes are evaluated empirically.}
  \label{tab:heuristic_modes}
  \begin{tabular}{lll}
    \toprule
    Mode & Definition & Admissibility \\
    \midrule
    \texttt{raw}               & $h_\theta(n)$                          & not guaranteed \\
    \texttt{scaled}            & $\alpha\,h_\theta(n)$                  & not guaranteed \\
    \texttt{clipped\_manhattan}& $\min(\alpha h_\theta(n), d_M(n))$     & guaranteed \\
    \bottomrule
  \end{tabular}
\end{table}

The admissibility audit on the validation split measures the fraction of cells
with $h_\theta(n)>h^\ast(n)+\epsilon$, the mean, 95th percentile, and maximum
overestimation in cell units, and how these quantities change with $\alpha$ and
with the clipped mode. The audit is paired with node-expansion measurements so
that the report can quantify the efficiency cost of enforcing admissibility.

\subsection{Search integration}
\texttt{NeuralHeuristic} loads the checkpoint once, builds the 29-dimensional
feature tensor for every cell of the maze once, and performs batched inference
under \texttt{torch.no\_grad()} with a default batch size of 8192. The
resulting cost field is cached as a two-dimensional array, so each heuristic
call during A$^\ast$ is an $O(1)$ lookup followed by the selected mode and
fallback logic. No feature construction or network forward pass occurs inside
the search loop.

The wrapper implements the same callable interface as
\texttt{ManhattanHeuristic} and \texttt{OracleHeuristic}, so the search code is
unchanged across heuristic families. This keeps the comparison controlled: the
same A$^\ast$ implementation, tie-breaking rule, and maze set are used for
Manhattan, oracle, and neural modes.

\subsection{Reproducibility}
The architecture diagram, model tests, and admissibility audit are reproduced
with:
\begin{verbatim}
python scripts/make_architecture_figure.py
python -m pytest tests/test_model.py tests/test_features.py \
    tests/test_heuristics.py tests/test_admissibility.py -q
python scripts/analyze_admissibility.py \
    --checkpoint models/neural_heuristic_seed42.pt
\end{verbatim}
The audit writes \texttt{results/admissibility/admissibility.csv},
\texttt{results/admissibility/admissibility\_summary.json}, and
\texttt{figures/admissibility\_cdf.png}. The checkpoint is loaded on CPU with
\texttt{torch.load(path, map\_location="cpu")}. Because recent PyTorch versions
default to \texttt{weights\_only=True}, the checkpoint contains only tensors and
primitive Python types (\texttt{list[float]} bin edges rather than NumPy
scalars), so it loads without disabling the safe loader.
